% Source: https://artofproblemsolving.com/community/c3h2470133_radicals
% Stable identity: aops-topic-2470133-q01
In $\sqrt[n]{b}$, the number under the radical sign is the \emph{radicand}, and the small number $n$ is the \emph{index} of the root (for example, in $\sqrt[3]{8}=2$ the index is $3$ and the radicand is $8$). ``Base'' and ``exponent'' are the names of $b$ and $n$ in a power $b^n$; note that $\sqrt[n]{b}=b^{\dfrac1n}$, where $b$ is the base and $\dfrac1n$ the exponent. The answer is $\boxed{\textbf{(D)}\ \text{radicand}}$.
